\section*{Abstract}

\singlespacing

More frequent gadget use in the classroom does not automatically make
learning more effective, so data are needed to examine its relationship
with students' learning motivation and social interaction. This
quantitative correlational study surveys \resN{} Grade 12 students of
Informatics II and III at SMAK Yos Sudarso Batam, selected by simple
random sampling, using a closed \resNItems-statement Likert questionnaire
analyzed in IBM SPSS 26.0 and independently reproduced with Python. Descriptive statistics summarize both
variables, and a scatter plot illustrates their relationship. Bivariate
analysis yields a strong positive Pearson correlation, $r = \resR$, with
a significance value of \resSig{} and an $R^2$ of \resRSquaredPct\%:
\resRSquaredPct\% of the
variation in motivation and social interaction is explained by gadget use
in the tested model while \resUnexplainedPct\% is not. The t-test
($t = \resT$) and
F-test ($F = \resF$) confirm statistical significance. Applying basic
probability to the same data, the probability of a high motivation score
is \resPB\%, rising to \resPBGivenA\% given a high gadget-use score. The
correlational design means these findings describe relationships, not
causation.

\noindent\textbf{Keywords:} gadget use, learning media, learning
motivation, social interaction, Pearson correlation, conditional
probability
